\section{Data}
The dataset used in this project is the UNSW-NB15 intrusion detection dataset \cite{UNSW}, which contains realistic network traffic data generated by the Cyber Range Lab of the Australian Centre for Cyber Security. It includes both normal network activity and multiple types of simulated cyber attacks such as DoS, exploits, fuzzers, worms, and reconnaissance. Each record represents a single network flow between two endpoints and is described by 49 features capturing different aspects of network behavior, including basic connection information (e.g., protocol, ports), temporal characteristics (e.g., duration, inter-arrival times), traffic volume (e.g., bytes and packets sent/received), and protocol-specific behavior.

\nextparagraph
The dataset contains both numerical and categorical attributes. Categorical features such as protocol type, service, and connection state are transformed into numerical representations using one-hot encoding. The target variable exists in two forms: a binary label indicating normal vs. attack traffic, and a multi-class label specifying the exact attack category. This allows the models to be evaluated both as binary intrusion detectors and as fine-grained attack classifiers.

\nextparagraph
A key property of the dataset is class imbalance, as certain attack types appear significantly more frequently than others. To address this issue, the Synthetic Minority Oversampling Technique (SMOTE) is applied to the training data, generating synthetic samples for underrepresented classes. This helps to reduce bias toward majority classes and improves the models’ ability to detect rare attack types.

\nextparagraph
Before training, all numerical features are normalized using Min-Max scaling to ensure that all features lie within a comparable range. This is particularly important for distance-based methods such as K-Means and for improving the convergence behavior of SVM-based models. The dataset does not contain missing values in its provided form, which simplifies preprocessing and ensures data consistency.

\nextparagraph
After preprocessing, the dataset is split into training and testing sets. The training set is further balanced using SMOTE for both binary and multi-class classification tasks, while the test set remains untouched to provide a realistic evaluation of model performance. The final processed datasets are stored in a structured format (Parquet files) to ensure efficient loading during model training and evaluation.